\documentclass[11pt, a4paper]{article}

% --- UNIVERSAL PREAMBLE BLOCK ---
\usepackage[a4paper, top=2.5cm, bottom=2.5cm, left=2cm, right=2cm]{geometry}
\usepackage{fontspec}

\usepackage[english, bidi=basic, provide=*]{babel}
\babelprovide[import, onchar=ids fonts]{english}

% Set default/Latin font to Sans Serif in the main (rm) slot for a modern look
\babelfont{rm}{Noto Sans}

% --- PACKAGES FOR VISUALS AND LAYOUT ---
\usepackage{xcolor}
\usepackage{tikz}
\usetikzlibrary{shapes.geometric, arrows.meta, positioning, shadows, calc, fit}
\usepackage[many]{tcolorbox}
\usepackage{multicol}
\usepackage{enumitem}

% Fix lists for non-English standard if needed (safe default)
\setlist[itemize]{label=\textcolor{primary}{\textbullet}}

% --- COLOR PALETTE ---
\definecolor{primary}{HTML}{1E3A8A}   % Deep Blue
\definecolor{secondary}{HTML}{0D9488} % Teal
\definecolor{accent}{HTML}{D97706}    % Amber
\definecolor{bglight}{HTML}{F8FAFC}   % Slate 50
\definecolor{textdark}{HTML}{0F172A}  % Slate 900
\definecolor{textmuted}{HTML}{64748B} % Slate 500

% --- CUSTOM TIKZ ICONS ---
% Since external icons (like fontawesome) are forbidden, we draw native vector icons
\newcommand{\iconGear}[1][1.2]{%
\begin{tikzpicture}[scale=#1, baseline=-0.3em]
  \fill[secondary] (0,0) circle (0.35);
  \fill[white] (0,0) circle (0.15);
  \foreach \i in {0, 45, ..., 315} {
    \fill[secondary, rotate=\i] (0.3, -0.08) rectangle (0.45, 0.08);
  }
\end{tikzpicture}%
}

\newcommand{\iconChart}[1][1.2]{%
\begin{tikzpicture}[scale=#1, baseline=-0.3em]
  \fill[primary] (-0.3,-0.3) rectangle (-0.1, 0.1);
  \fill[primary] (0,-0.3) rectangle (0.2, 0.4);
  \fill[accent] (0.3,-0.3) rectangle (0.5, 0.6);
  \draw[textdark, thick, line cap=round] (-0.4,-0.4) -- (0.6,-0.4);
  \draw[textdark, thick, line cap=round] (-0.4,-0.4) -- (-0.4,0.7);
\end{tikzpicture}%
}

\newcommand{\iconBulb}[1][1.2]{%
\begin{tikzpicture}[scale=#1, baseline=-0.3em]
  \fill[accent] (0,0.2) circle (0.25);
  \fill[accent] (-0.15,0) rectangle (0.15,0.2);
  \fill[textdark] (-0.1,-0.1) rectangle (0.1,0.05);
  \fill[textdark] (-0.05,-0.2) rectangle (0.05,-0.1);
\end{tikzpicture}%
}

% --- TCOLORBOX STYLES ---
\tcbset{
  % Style for the main executive summary
  summarybox/.style={
    enhanced,
    colback=bglight,
    colframe=primary,
    boxrule=0pt,
    leftrule=4pt,
    arc=2pt,
    fonttitle=\bfseries\Large,
    coltitle=primary,
    colbacktitle=bglight,
    title=#1,
    attach title to upper,
    after title={\\[0.5em]},
    drop shadow=black!10
  },
  % Style for detailed component features
  featurebox/.style={
    enhanced,
    colback=white,
    colframe=secondary,
    boxrule=1pt,
    arc=4pt,
    fonttitle=\bfseries,
    title=#1,
    drop shadow=black!5,
    toptitle=1.5mm,
    bottomtitle=1.5mm
  }
}

% Hyperref must be the last package
\usepackage[colorlinks=true, linkcolor=primary, urlcolor=secondary]{hyperref}

\begin{document}

% --- HEADER ---
\begin{center}
    \begin{tikzpicture}
        % Background header ribbon
        \fill[primary, rounded corners=3pt] (-8.5, -1.2) rectangle (8.5, 1.2);
        \fill[secondary, rounded corners=3pt] (-8.5, -1.2) rectangle (-7.5, 1.2);
        
        % Title text
        \node[text=white, font=\Huge\bfseries] at (0, 0.2) {Project Nova Overview};
        \node[text=white!80, font=\large] at (0, -0.6) {Next-Generation Analytics Infrastructure};
    \end{tikzpicture}
\end{center}

\vspace{0.5cm}

% --- EXECUTIVE SUMMARY ---
\begin{tcolorbox}[summarybox={\iconBulb[1.5] Executive Summary}]
    \color{textdark}
    Project Nova is a highly scalable, real-time analytics engine designed to process millions of data points per second. By leveraging cutting-edge container orchestration and distributed databases, it provides unparalleled visibility into user behavior and system performance. 
    
    This document outlines the core architecture, detailed component breakdown, and expected performance metrics of the proposed system.
\end{tcolorbox}

\vspace{0.8cm}

% --- SYSTEM ARCHITECTURE (VISUALIZATION) ---
\noindent\textcolor{primary}{\Large\bfseries \iconGear[1.5] Core Architecture Visualization}
\vspace{0.3cm}

\begin{center}
\begin{tikzpicture}[
    node distance=1.5cm and 2cm,
    box/.style={rectangle, draw=primary, thick, fill=white, minimum width=3cm, minimum height=1.5cm, rounded corners=4pt, drop shadow},
    db/.style={cylinder, draw=accent, thick, fill=white, shape border rotate=90, aspect=0.25, minimum width=2.5cm, minimum height=2.5cm, drop shadow},
    arrow/.style={-{Stealth[scale=1.2]}, thick, draw=textmuted}
]

    % Nodes
    \node[box] (client) {\begin{tabular}{c}\textbf{Client App}\\(Web / Mobile)\end{tabular}};
    \node[box, right=of client, yshift=2cm, fill=bglight] (api1) {\begin{tabular}{c}\textbf{Ingestion API}\\(Node.js)\end{tabular}};
    \node[box, right=of client, yshift=-2cm, fill=bglight] (api2) {\begin{tabular}{c}\textbf{Query API}\\(Python/FastAPI)\end{tabular}};
    
    \node[box, right=of api1, fill=secondary!10, draw=secondary] (kafka) {\begin{tabular}{c}\textbf{Message Stream}\\(Kafka)\end{tabular}};
    
    \node[db, right=of api2, xshift=1.5cm] (db) {\begin{tabular}{c}\textbf{Data Lake}\\(ClickHouse)\end{tabular}};
    
    % Grouping box
    \node[draw=textmuted, dashed, rounded corners, fit=(api1) (api2) (kafka), inner sep=0.5cm, label=above:\textcolor{textmuted}{Microservices Layer}] (layer) {};

    % Connections
    \draw[arrow] (client) -- node[above, sloped, text=textmuted, font=\small] {Events} (api1.west);
    \draw[arrow, <->] (client) -- node[above, sloped, text=textmuted, font=\small] {GraphQL} (api2.west);
    
    \draw[arrow] (api1) -- (kafka);
    \draw[arrow] (kafka) -| node[right, near start, text=textmuted, font=\small] {Batch load} (db);
    \draw[arrow, <->] (api2) -- node[above, text=textmuted, font=\small] {SQL} (db);

\end{tikzpicture}
\end{center}

\vspace{0.8cm}

% --- DETAILED COMPONENTS ---
\noindent\textcolor{primary}{\Large\bfseries Detailed Component Analysis}
\vspace{0.3cm}

\begin{multicols}{2}

\begin{tcolorbox}[featurebox={1. High-Throughput Ingestion}]
    The \textbf{Ingestion API} is optimized for write-heavy workloads. It validates incoming telemetry data and acts as a buffer.
    \begin{itemize}[leftmargin=*]
        \item Sub-millisecond latency.
        \item Horizontally scalable via Kubernetes.
        \item Automatic schema validation.
    \end{itemize}
\end{tcolorbox}

\begin{tcolorbox}[featurebox={2. Distributed Stream Processing}]
    Using \textbf{Kafka}, we decouple data ingestion from storage. This allows multiple consumers to process the same data streams in real-time.
    \begin{itemize}[leftmargin=*]
        \item Fault-tolerant architecture.
        \item 7-day data retention buffer.
    \end{itemize}
\end{tcolorbox}

\columnbreak

\begin{tcolorbox}[featurebox={3. Columnar Data Lake}]
    \textbf{ClickHouse} serves as the analytical backend, executing aggregations over billions of rows in milliseconds.
    \begin{itemize}[leftmargin=*]
        \item 10x compression ratio.
        \item Materialized views for fast dashboards.
        \item Vectorized query execution.
    \end{itemize}
\end{tcolorbox}

\begin{tcolorbox}[featurebox={4. Unified Query Layer}]
    A \textbf{GraphQL API} abstracts the underlying database complexity, providing frontend engineers with a simple interface.
    \begin{itemize}[leftmargin=*]
        \item Role-based access control (RBAC).
        \item Query caching via Redis.
    \end{itemize}
\end{tcolorbox}

\end{multicols}

\vspace{0.5cm}

% --- DATA VISUALIZATION / METRICS ---
\noindent\textcolor{primary}{\Large\bfseries \iconChart[1.5] Projected Performance Gains}
\vspace{0.3cm}

We expect significant improvements in both ingestion speed and query latency compared to the legacy system. The chart below illustrates the targeted benchmarks.

\vspace{0.4cm}
\begin{center}
\begin{tikzpicture}[font=\sffamily]
    % Grid lines
    \draw[gray!30, dashed] (0, 0) grid (10, 5);
    
    % Axes
    \draw[thick, textdark, -{Stealth}] (0,0) -- (10.5,0) node[right] {Metric};
    \draw[thick, textdark, -{Stealth}] (0,0) -- (0,5.5) node[above] {Relative Performance (x)};
    
    % Y-axis labels
    \foreach \y in {1,2,3,4,5}
        \node[left, textmuted] at (0,\y) {\y x};
        
    % Bars
    % Group 1: Query Speed
    \fill[secondary] (1.5,0) rectangle (2.5, 1);
    \node[above] at (2, 1) {Legacy};
    
    \fill[primary] (2.8,0) rectangle (3.8, 4.5);
    \node[above, font=\bfseries] at (3.3, 4.5) {4.5x};
    
    \node[below, font=\bfseries, textdark] at (2.65, -0.2) {Query Speed};

    % Group 2: Throughput
    \fill[secondary] (6.5,0) rectangle (7.5, 1);
    \node[above] at (7, 1) {Legacy};
    
    \fill[accent] (7.8,0) rectangle (8.8, 3.2);
    \node[above, font=\bfseries] at (8.3, 3.2) {3.2x};
    
    \node[below, font=\bfseries, textdark] at (7.65, -0.2) {Data Throughput};

    % Legend
    \matrix [draw=textmuted, rounded corners, fill=white, below left] at (10, 5) {
      \fill[secondary] (0,0) rectangle (0.4,0.3); & \node {Legacy System}; \\
      \fill[primary] (0,0) rectangle (0.4,0.3); & \node {Nova (Read)}; \\
      \fill[accent] (0,0) rectangle (0.4,0.3); & \node {Nova (Write)}; \\
    };

\end{tikzpicture}
\end{center}

\end{document}